% ------------------------------------------------------------------------
% PAKIETY – WERSJA STABILNA (RAFMAT / WINDOWS SERVER)
% ------------------------------------------------------------------------

% --- Kodowanie i język ---
\usepackage[utf8]{inputenc} % To jest standard dla UTF-8
\usepackage[T1]{fontenc}    % Poprawne wyświetlanie polskich znaków w PDF
\usepackage[polish]{babel}
\usepackage{csquotes}

% --- Matematyka (zostawione, nie przeszkadzają) ---
\usepackage{amsmath}
\usepackage{amssymb}
\usepackage{mathrsfs}

% --- Grafika i układ strony ---
\usepackage{graphicx}
\usepackage{float}
\usepackage{geometry}
\geometry{
  a4paper,
  left=3cm,
  right=2.5cm,
  top=3cm,
  bottom=3cm
}

% --- Kolory i listingi ---
\usepackage{xcolor}
\usepackage{listings}
\usepackage[normalem]{ulem} % <--- DODANO (do podkreślania linków w pkt 11)

% --- Bibliografia ---
\usepackage[
  backend=biber,
  style=numeric,
  sorting=none
]{biblatex}
\addbibresource{bibliografia.bib}

% --- Linki (ZAWSZE po biblatex!) ---

\usepackage{xurl}      % ← MUSI BYĆ WCZEŚNIEJ (lamanie linkow)
\usepackage[
  colorlinks=true,
  urlcolor=blue,
  linkcolor=black,
  citecolor=blue,
  breaklinks=true,
  draft=false,   % ← linki zawsze niebieskie
  bookmarksopen=true   % linkowanie zdjec
]{hyperref}

% --- Spisy i sekcje ---
\usepackage{tocloft}
\usepackage{titlesec}
\usepackage{enumitem} % -> do ściskania list - żeby się na 1 stronie zmieścił rozdział 1

% --- Strona tytułowa (uczelniany pakiet) ---
\usepackage{strona_tytulowa}

% --- TODO (do wyłączenia przed oddaniem) ---
\usepackage[hide]{todo}

% --- Tabele ---
\usepackage{array}
\usepackage{longtable}
\usepackage{booktabs}

% --- Nagłówki i stopki ---
\usepackage{fancyhdr}
\setlength{\headheight}{24pt}

% --- Nagłówki i stopki ---
\pdfminorversion=7

% --- Listing styl ---
\definecolor{codegreen}{rgb}{0,0.6,0}
\definecolor{codegray}{rgb}{0.5,0.5,0.5}
\definecolor{codepurple}{rgb}{0.58,0,0.82}
\definecolor{backcolour}{rgb}{0.95,0.95,0.92}

\lstdefinestyle{mystyle}{
	backgroundcolor=\color{backcolour},
	commentstyle=\color{codegreen},
	keywordstyle=\color{blue},
	numberstyle=\tiny\color{codegray},
	stringstyle=\color{codepurple},
	basicstyle=\ttfamily\footnotesize,
	breaklines=true,
	numbers=left,
	numbersep=5pt,
	showstringspaces=false,
	tabsize=2
}

\lstset{style=mystyle}



% ============================
% Definicja języka PowerShell dla listings
% ============================
\lstdefinelanguage{PowerShell}{
  morekeywords={
    Add-ADUser, New-ADUser, New-ADGroup, New-ADOrganizationalUnit,
    Get-ADUser, Get-ADGroup, Import-Csv, ForEach-Object,
    If, Else, Foreach, While, Function, Param
  },
  sensitive=true,
  morecomment=[l]{\#},
  morestring=[b]"
}


\usepackage{placeins}
